\scanpage{42}
Lorsqu'on se fixe un point $\xi_0\in S(\mathbb{C})$,\deleted{\unsure{$\xi$|,}}\inserted{et qu'on prend $T=T_{\xi_0}$,} alors la donnée du revêtement principal $\bar S_{\mathbb{C}}$ et de $\deleted{\unsure{f|p}}p$ équivaut à la donnée de
\[
\begin{aligned}
&\deleted{\unsure{w(\pi_1(S(\mathbb{C}),\xi_0)|w:\pi_1(S(\mathbb{C}),\xi_0)}}\\
&\unsure{w|\omega}:\pi_1(S(\mathbb{C}),\xi_0)\longrightarrow G(\mathbb{Q})\subset G(\mathbb{C})
\end{aligned}
\]
[il est associé à la connexion intégrable de $P$, qui en définit une sur $P_{\mathbb{C}}^{\mathrm{an}}$ sur $S_{\mathbb{C}}^{\mathrm{an}}$, d'où $\Gamma=\pi_1(S(\mathbb{C}),\xi_0)\longrightarrow G(\mathbb{C})$, dont il faut demander qu'il se factorise par $G(\mathbb{Q})$] et à un \deleted{hom du revêtement universel} point $\alpha_0$ de $P(\xi_0)$ i.e. une trivialisation de $P_{\xi_0}$, qui définit alors un unique morphisme \inserted{$q$} du rev. universel $\bar S_{\mathbb{C}}$ de $S_{\mathbb{C}}$ dans $P_{\mathbb{C}}^{\mathrm{an}}=P(\mathbb{C})$, appliquant le pt marqué $\bar\xi_0$ en $\alpha_0$ ; il satisfait \deleted{\unsure{à|la}} $q(\bar s\gamma)=q(\bar s)\unsure{w|W}(\gamma)$ pour tout $\hat s\in\bar S_{\mathbb{C}}$, \deleted{$\unsure{\gamma\in\pi_1(S|\gamma\in\pi_1}$} $\gamma\in\Pi_1(S(\mathbb{C}),\xi_0)$, \deleted{Le morphisme \unsure{satisfai|satisfait}} et définit
\[
\widetilde P_{\mathbb{C}}=\overline{\unsure{S|\mathbb{S}}}_{\mathbb{C}}\mathbin{\times}^{\Gamma}G(\mathbb{Q})\hookrightarrow P_{\mathbb{C}}^{\mathrm{an}}
\]
[\unsure{restriction|restiction} du groupe structural \inserted{du $P_{\mathbb{C}}^{\mathrm{an}}$} de $G(\mathbb{C})$ à $G(\mathbb{Q})$].

Comme structure importante, il y a une application
\[
f:\widetilde P_{\mathbb{C}}\longrightarrow\unsure{\gamma|\mathscr{Y}}_G
\]
où $\gamma_G$ est l'ens. des structures de Hodge sur $G_{\mathbb{R}}$ (\deleted{associées,}\inserted{compatibles,} si on y tient, avec un morphisme $i:\mathbb{G}_{m\,\mathbb{Q}}\longrightarrowG$ donné). Cette application $f$ satisfait les conditions

a) $f$ holomorphe

b) $f(xg)=g^{-1}f(x)$ \qquad $g\in G(\mathbb{Q}),\,x\in\widetilde P_{\mathbb{C}}$

\unsure{b|c}) $\forall x\in\widetilde P_{\mathbb{C}}$ sur le point $s$ de $S$, la \inserted{structure de} filtration sur $T_{\mathbb{C}}$ sur $\operatorname{Rep}(\unsure{G|C})$, définie par $f(\unsure{x|s}):\mathbb{S}\longrightarrow G_{\mathbb{R}}$, est compatible
% archive sheet number: 41
